%%%%%%%%%%%%%%%%%%%%%%%%
%
% $Autor: Hemanth Jadiswami Prabhakaran $
% $Datum: 2026-09-25 $
% $Pfad: GitHub/MBIDA_Master_Thesis/Manual/Chapters/en/05Functions.tex $
% $Version: 1 $
%
% $Project: MBIDA Thesis - Manual $
%
%%%%%%%%%%%%%%%%%%%%%%%%


\chapter{Functions}
\label{chap:functions}

This chapter describes each function of the dashboard: what it does, how it is used and what to look out for. The functions fall into three groups -- \emph{choosing what to look at} (Sections~\ref{sec:fnHierarchy}--\ref{sec:fnSearch}), \emph{choosing which forecasts to compare} (Sections~\ref{sec:fnBase}--\ref{sec:fnRec}) and \emph{exploring the chart} (Sections~\ref{sec:fnChart}--\ref{sec:fnTheme}). Section~\ref{sec:fnTasks} closes with typical analysis tasks that combine them.

\section{The Sales Hierarchy}
\label{sec:fnHierarchy}

All navigation in the dashboard follows the six-level hierarchy in Figure~\ref{fig:hierarchyTree}. Every node is the sum of the nodes directly below it: the three states add up to the total, the four Californian stores add up to \VAL{CA}, and so on down to the 30,240 single products-in-a-store at the bottom.

\begin{figure}[H]
  \centering
  \resizebox{\linewidth}{!}{% The six-level sales hierarchy. The highlighted path is the example used
% throughout the manual; grey nodes stand for siblings that are not expanded.
\begin{tikzpicture}[
    n/.style={draw=ManBlue,fill=ManBlueLight,rounded corners=2pt,font=\ttfamily\scriptsize,
              minimum height=0.55cm,inner sep=3pt},
    on/.style={n,draw=AppRed,fill=AppRed!12,very thick},
    off/.style={n,draw=ManGrey!60,fill=ManGreyLight,text=ManGrey},
    e/.style={draw=ManGrey!70,thin},
    eon/.style={draw=AppRed,thick},
    lvl/.style={anchor=east,font=\sffamily\small\bfseries,text=ManBlue},
    cnt/.style={anchor=west,font=\sffamily\small,text=ManGrey!70!black},
  ]
  \def\dy{1.35}

  % level labels (left) and node counts (right)
  \foreach \i/\L/\c in {0/Total/1, 1/State/3, 2/Store/10, 3/Category/30,
                        4/Department/70, 5/Item/{30,240}}{
     \pgfmathtruncatemacro{\lv}{\i+1}\node[lvl] at (-0.6,-\i*\dy) {\lv\ \ \L};
     \node[cnt] at (12.4,-\i*\dy) {\c\ series};
  }
  \draw[ManGrey!40] (12.3,0.4) -- (12.3,-5.4*\dy);

  % level 1
  \node[on] (T) at (5.8,0) {Total};
  % level 2
  \node[on]  (CA) at (2.8,-\dy) {CA};
  \node[off] (TX) at (5.8,-\dy) {TX};
  \node[off] (WI) at (8.8,-\dy) {WI};
  % level 3
  \node[on]  (C1) at (0.6,-2*\dy) {CA\_1};
  \node[off] (C2) at (2.0,-2*\dy) {CA\_2};
  \node[off] (C3) at (3.4,-2*\dy) {CA\_3};
  \node[off] (C4) at (4.8,-2*\dy) {CA\_4};
  \node[off] (TXs) at (6.6,-2*\dy) {TX\_1..3};
  \node[off] (WIs) at (9.2,-2*\dy) {WI\_1..3};
  % level 4
  \node[on]  (F)  at (0.6,-3*\dy) {FOODS};
  \node[off] (HB) at (2.5,-3*\dy) {HOBBIES};
  \node[off] (HH) at (4.6,-3*\dy) {HOUSEHOLD};
  % level 5
  \node[off] (F1) at (0.1,-4*\dy) {FOODS\_1};
  \node[off] (F2) at (1.85,-4*\dy) {FOODS\_2};
  \node[on]  (F3) at (3.6,-4*\dy) {FOODS\_3};
  % level 6
  \node[off] (I1) at (1.6,-5*\dy) {FOODS\_3\_001};
  \node[font=\sffamily\small,text=ManGrey] (dots) at (3.35,-5*\dy) {\dots};
  \node[on]  (I2) at (5.1,-5*\dy) {FOODS\_3\_090};
  \node[font=\sffamily\small,text=ManGrey] (dots2) at (6.8,-5*\dy) {\dots};

  % edges
  \draw[eon] (T.south) -- (CA.north);
  \draw[e] (T.south) -- (TX.north);  \draw[e] (T.south) -- (WI.north);
  \draw[eon] (CA.south) -- (C1.north);
  \foreach \c in {C2,C3,C4} \draw[e] (CA.south) -- (\c.north);
  \draw[e] (TX.south) -- (TXs.north); \draw[e] (WI.south) -- (WIs.north);
  \draw[eon] (C1.south) -- (F.north);
  \draw[e] (C1.south) -- (HB.north); \draw[e] (C1.south) -- (HH.north);
  \draw[e] (F.south) -- (F1.north); \draw[e] (F.south) -- (F2.north);
  \draw[eon] (F.south) -- (F3.north);
  \draw[e] (F3.south) -- (I1.north);
  \draw[eon] (F3.south) -- (I2.north);

  % node label of the highlighted path
  \node[anchor=west,draw=AppRed,fill=white,rounded corners=2pt,font=\ttfamily\scriptsize,
        align=left,inner sep=4pt] at (6.9,-3.55*\dy)
        {\sffamily\bfseries Node label shown in the app:\\[1pt]
         Total/CA/CA\_1/FOODS/FOODS\_3/\\FOODS\_3\_090 \ (level = 6)};

  % "sums up" arrow
  \draw[-{Stealth[length=2.5mm]},very thick,ManGreen] (11.5,-5*\dy) -- (11.5,0)
       node[midway,right,font=\sffamily\scriptsize,text=ManGreen,align=left,xshift=1pt]{sums\\up};
\end{tikzpicture}
}
  \caption{The six-level sales hierarchy with the example path used in this manual}
  \label{fig:hierarchyTree}
\end{figure}

\begin{table}[H]
\centering
\caption{Levels of the hierarchy}
\label{tab:fnLevels}
\small
\begin{tabularx}{\linewidth}{@{}clrX@{}}
\toprule
\textbf{Level} & \textbf{Name} & \textbf{Nodes} & \textbf{Example} \\
\midrule
1 & Total      & 1      & all sales \\
2 & State      & 3      & \VAL{CA} (California), \VAL{TX} (Texas), \VAL{WI} (Wisconsin) \\
3 & Store      & 10     & \VAL{CA\_1} \dots\ \VAL{CA\_4}, \VAL{TX\_1} \dots\ \VAL{TX\_3}, \VAL{WI\_1} \dots\ \VAL{WI\_3} \\
4 & Category   & 30     & \VAL{FOODS} in store \VAL{CA\_1} \\
5 & Department & 70     & \VAL{FOODS\_3} in store \VAL{CA\_1} \\
6 & Item       & 30,240 & product \VAL{FOODS\_3\_090} in store \VAL{CA\_1} \\
\midrule
  & \textbf{Sum} & \textbf{30,354} & \\
\bottomrule
\end{tabularx}
\end{table}

\begin{NoteBox}[Categories are always inside a store]
The hierarchy is strictly \emph{store first, then category}. The dashboard therefore shows ``\VAL{FOODS} in store \VAL{CA\_1}'', but not ``\VAL{FOODS} in all stores together''. Such a cross-store view is not one of the 30,354 nodes (Section~\ref{sec:specLimits}).
\end{NoteBox}

\section{Navigating with the Cascading Dropdowns}
\label{sec:fnNavigate}

\paragraph{What it does.} The five dropdowns \Callout{2} select the node. The dashboard reads them from top to bottom and stops at the first one that is left on \VAL{(All)}; the node is everything chosen above that point (Figure~\ref{fig:cascadeLogic}).

\begin{figure}[H]
  \centering
  % How the five dropdowns are turned into the node that is plotted.
% Reading stops at the first level left on "(All)".
\begin{tikzpicture}[
    q/.style={man/decision,text width=2.3cm,font=\sffamily\scriptsize},
    r/.style={man/boxgreen,text width=4.6cm,font=\sffamily\scriptsize,minimum height=0.8cm,align=left},
    node distance=0.55cm,
  ]
  \node[man/terminal,font=\sffamily\scriptsize] (s) {Read the sidebar from the top};
  \node[q,below=of s] (q1) {State\\= (All)?};
  \node[q,below=of q1] (q2) {Store\\= (All)?};
  \node[q,below=of q2] (q3) {Category\\= (All)?};
  \node[q,below=of q3] (q4) {Department\\= (All)?};
  \node[q,below=of q4] (q5) {Item\\= (All)?};
  \node[r,below=of q5,man/boxorange,text width=9.5cm,align=center] (r6) {\textbf{Level 6}\; one item in one store\\ \texttt{Total/CA/CA\_1/FOODS/FOODS\_3/FOODS\_3\_090}};

  \node[r,right=1.2cm of q1] (r1) {\textbf{Level 1}\; grand total\; \texttt{Total}};
  \node[r,right=1.2cm of q2] (r2) {\textbf{Level 2}\; one state\; \texttt{Total/CA}};
  \node[r,right=1.2cm of q3] (r3) {\textbf{Level 3}\; one store\; \texttt{Total/CA/CA\_1}};
  \node[r,right=1.2cm of q4] (r4) {\textbf{Level 4}\; one category in a store\\ \texttt{Total/CA/CA\_1/FOODS}};
  \node[r,right=1.2cm of q5] (r5) {\textbf{Level 5}\; one department in a store\\ \texttt{Total/CA/CA\_1/FOODS/FOODS\_3}};

  \draw[man/arrow] (s) -- (q1);
  \foreach \a/\b in {q1/q2,q2/q3,q3/q4,q4/q5,q5/r6}
     \draw[man/arrow] (\a) -- node[man/label,right]{no} (\b);
  \foreach \i in {1,...,5}
     \draw[man/arrow] (q\i) -- node[man/label,above]{yes} (r\i);

\end{tikzpicture}

  \caption{How the five dropdowns determine the node that is shown}
  \label{fig:cascadeLogic}
\end{figure}

\paragraph{How to use it.}
\begin{itemize}
  \item \textbf{Going down:} choose values from top to bottom. Each choice unlocks the next dropdown and restricts it to what exists below.
  \item \textbf{Going up one level:} set the lowest chosen dropdown back to \VAL{(All)}.
  \item \textbf{Going back to the total:} set \UI{State} to \VAL{(All)}; everything below follows.
  \item \textbf{Going sideways:} change a dropdown to a sibling, e.g.\ \UI{Store} from \VAL{CA\_1} to \VAL{CA\_2}. The dropdowns below it are usually reset to \VAL{(All)}, because their lists of values change with the new store.
\end{itemize}

\paragraph{Locked dropdowns.} A dropdown below an \VAL{(All)} offers only \VAL{(All)} (Figure~\ref{fig:fnLocked}). This is intended: without a store, for example, ``department \VAL{FOODS\_3}'' would be ambiguous.

\begin{figure}[H]
  \centering
  \Screenshot[0.45\linewidth]{Images/05Functions/LockedDropdowns.png}
  \caption{With \UI{State} on \VAL{(All)}, all lower dropdowns are locked to \VAL{(All)}}
  \label{fig:fnLocked}
\end{figure}

\section{Searching in a Dropdown}
\label{sec:fnSearch}

Every dropdown can be searched. Click it and start typing; the list is filtered to entries containing the typed text anywhere in the name (Figure~\ref{fig:fnSearch}). Press \keys{Enter} to take the first match or click the wanted entry. This is the fastest way to reach an item: a department contains up to several hundred items per store, which is too many to scroll through.

\begin{figure}[H]
  \centering
  \Screenshot[0.5\linewidth]{Images/05Functions/DropdownSearch.png}
  \caption{Typing \VAL{090} in the \UI{Item} dropdown}
  \label{fig:fnSearch}
\end{figure}

\begin{TipBox}
Item names follow the pattern \texttt{CATEGORY\_DEPARTMENT\_NUMBER}, e.g.\ \VAL{HOUSEHOLD\_1\_118}. Typing just the number is usually enough, because the department is already fixed by the dropdown above.
\end{TipBox}

\section{Choosing the Base Model}
\label{sec:fnBase}

\paragraph{What it does.} The \UI{Base model} buttons \Callout{3} select which of three forecasting models draws the dotted line -- and which model the reconciled lines are based on. All three were calculated in advance for every node, so switching is instant.

\begin{table}[H]
\centering
\caption{The three base models}
\label{tab:fnModels}
\small
\begin{tabularx}{\linewidth}{@{}lXX@{}}
\toprule
\textbf{Button} & \textbf{Idea in one sentence} & \textbf{Typical look in the chart} \\
\midrule
\VAL{histavg} & The forecast is the average of all past months. & A flat line; serves as the simple benchmark every other model must beat. \\
\VAL{ets} & Exponential smoothing: recent months count more, and trend and yearly season are modelled \cite{Hyndman:2021}. & Follows level, trend and the December peak. \\
\VAL{theta} & Combines a long-term trend line with a short-term smoothed view \cite{Hyndman:2021}. & Similar to \VAL{ets}, often slightly smoother. \\
\bottomrule
\end{tabularx}
\end{table}

\paragraph{What to look out for.} For nodes with too little history to fit \VAL{ets} or \VAL{theta} -- mostly single items that rarely sell -- the forecast falls back to the historical average. Such a node shows the same flat line for all three buttons. This is expected and is not an error.

\section{Choosing Reconciliation Methods}
\label{sec:fnRec}

\paragraph{What it does.} Base forecasts are made for each node separately, so the forecasts of the parts do not add up to the forecast of the whole. Reconciliation corrects this. The \UI{Reconciliation method(s)} box \Callout{4} adds one solid line per selected method.

\begin{table}[H]
\centering
\caption{The three reconciliation methods}
\label{tab:fnRec}
\small
\begin{tabularx}{\linewidth}{@{}lX@{}}
\toprule
\textbf{Method} & \textbf{Idea in plain words} \\
\midrule
\VAL{Bottom-Up} & Trust the item forecasts and add them up to every higher level. \\
\VAL{Top-Down} & Trust the total forecast and split it downwards according to each node's historical share of sales. \\
\VAL{MinTrace (wls\_struct)} & A compromise: adjust the forecasts on all levels at once, as little as possible, so that they add up; nodes that consist of many items get more weight \cite{Wickramasuriya:2019}. \\
\bottomrule
\end{tabularx}
\end{table}

\paragraph{How to use it.} Any combination is possible: none (only actual and base line), one, two or all three methods. Comparing all three at once is the main use. Because each reconciled line is the \emph{base forecast after adjustment}, the gap between the dotted line and a solid line shows how much the method changed the forecast for this node.

\paragraph{Reading the effect by level.}
\begin{itemize}
  \item At the \textbf{Total} (level~1), \VAL{Top-Down} is identical to the base forecast, because it starts from the total.
  \item At the \textbf{Item} level (level~6), \VAL{Bottom-Up} is identical to the base forecast, because it starts from the items.
  \item \VAL{MinTrace} changes every level a little.
\end{itemize}

\section{Exploring the Chart}
\label{sec:fnChart}

The chart is interactive \cite{Plotly:2026}. Table~\ref{tab:fnChart} lists all interactions.

\begin{table}[H]
\centering
\caption{Chart interactions}
\label{tab:fnChart}
\small
\begin{tabularx}{\linewidth}{@{}lX@{}}
\toprule
\textbf{Action} & \textbf{Result} \\
\midrule
Mouse over a line & Box with month and value of that point \\
Drag a rectangle  & Zoom into this area \\
Drag along an axis & Zoom or move only this axis \\
Double-click      & Back to the full view \\
Click a legend entry & Hide or show this line \\
Double-click a legend entry & Show only this line; double-click again shows all \\
Camera button \Callout{10} & Save the current view as \ac{png} \\
\bottomrule
\end{tabularx}
\end{table}

\begin{figure}[H]
  \centering
  \Screenshot[0.49\linewidth]{Images/05Functions/HoverTooltip.png}
  \hfill
  \Screenshot[0.49\linewidth]{Images/05Functions/ZoomedTestPeriod.png}
  \caption{Hover information (left) and the test period after zooming in (right)}
  \label{fig:fnChart}
\end{figure}

\begin{NoteBox}
Zoom and hidden lines are reset whenever another node, model or method is selected, because the chart is drawn anew. Save an image first if the view is needed later.
\end{NoteBox}

\section{Light and Dark Mode}
\label{sec:fnTheme}

The dashboard follows the theme chosen in the page menu \Callout{11} (Section~\ref{sec:uiMenu}). The chart detects the active theme and switches its colours for lines, grid, text and shading, so it stays readable in both modes. With \VAL{Use system setting}, the theme changes together with the operating system's light/dark setting. No separate switch inside the dashboard is needed.

\section{Typical Analysis Tasks}
\label{sec:fnTasks}

\begin{table}[H]
\centering
\caption{Typical tasks and how to do them}
\label{tab:fnTasks}
\small
\begin{tabularx}{\linewidth}{@{}p{4.3cm}X@{}}
\toprule
\textbf{Task} & \textbf{Procedure} \\
\midrule
How good is the forecast for my store? & Select the store, keep \VAL{MinTrace}, try all three base models; compare with the black line in the \UI{Test} region. \\
Which reconciler suits a department? & Select the department, one base model, all three methods; hover along the test year. \\
How much does reconciliation change a forecast? & Compare the dotted base line with the solid lines; the larger the gap, the stronger the adjustment. \\
Why are item forecasts worse than totals? & Compare the total (level~1) with any item (level~6): item sales are small and irregular, totals are smooth. \\
Put a chart into a report & Set up the view, zoom if needed, click the camera button; switch to light mode first for printing. \\
\bottomrule
\end{tabularx}
\end{table}

The dashboard shows forecasts visually. The measured accuracy of all models and methods over the test year, per level, is summarised in Section~\ref{sec:specAccuracy}.
